\section{Implementation details}\label{app:impl}
\subsection{Data splits}
\subsubsection{Network level}
Days are UTC; the unit is the host-hour of Zeek or NFStream flows.
\begin{table}[!htbp]\centering\footnotesize\setlength{\tabcolsep}{3pt}
\caption{Train / calibration / test days for the network-level experiments. Train and calibration days contain no attack flow.}
\begin{tabularx}{\textwidth}{@{}lYYY@{}}
\toprule
Dataset & Train & Calibration & Test \\
\midrule
DEDALE (green zone) & 2024-12-24 -- 2025-01-02 & 2025-01-03 -- 2025-01-05 & 2025-01-06 -- 2025-01-19 \\
Unraveled & 2021-05-28 -- 2021-06-04 & 2021-06-05 -- 2021-06-07 & 2021-06-08 -- 2021-07-06 \\
OpTC (Zeek) & 2019-09-16 -- 2019-09-19 & 2019-09-20 -- 2019-09-22 & 2019-09-23 -- 2019-09-25 \\
\bottomrule
\end{tabularx}
\end{table}
\subsubsection{Provenance level}
The splits are PIDSMaker's defaults for each dataset.
\begin{table}[!htbp]\centering\footnotesize\setlength{\tabcolsep}{3pt}
\caption{Train / validation / test dates for the host-level experiments: first and last day, and the number of days actually used in parentheses (the days of a split need not be contiguous).}
\begin{tabularx}{\textwidth}{@{}lYYY@{}}
\toprule
Dataset & Train & Validation & Test \\
\midrule
%%SPLITS_PROV%%
\bottomrule
\end{tabularx}
\end{table}
LANL: development = days 2--12, test = days 13--30 of the red-team events (day index from the start of the data).

\subsection{Models and hyperparameters}
\begin{itemize}
\item Network features (per host-hour): flows out, distinct destination IPs and ports, share of failed
connections, new and rare external destinations (rare = contacted by at most two internal hosts before), beacon
regularity over 72\,h, bytes out and in, new admin-port edges, GNN edge surprise. Counts and bytes are transformed
with $\log(1+x)$; ratios are left unchanged. ``New'' means not seen on any earlier day.
\item Per-host tail model: one-sided empirical upper-tail probability of each feature against the host's own
training hours (all hosts' hours if the host has fewer than 24); channel score = most extreme feature (+0.01 $\times$
channel mean); conformal $p$-values from cross-fitted calibration (training days split by parity, plus the calibration
days); hour score $=-\log_{10}(\min_c p_c \times 6)$.
\item Pooled baselines: ECOD (pyod defaults); Isolation Forest (200 trees); autoencoder = scikit-learn
\texttt{MLPRegressor} with hidden layers (32, 8, 32), 300 iterations, on standardised features. Seed 42 unless stated.
\item GNN lateral channel: learnable 32-d node embeddings, two GraphSAGE-mean layers, MLP decoder over
(source, destination, port embedding of size 8); 300 epochs of Adam (learning rate 0.01, weight decay $10^{-4}$);
three random destinations and two random ports per observed triple as negatives; training days scored out of sample
by two models trained on alternate days; seed 42.
\item Velox / Velox-lite: PIDSMaker's Velox (tuned configuration, embedding size 32) and our re-implementation
(word2vec embeddings of node labels with PIDSMaker's tokenizer, size 32, 10 epochs; a linear layer to 128 dimensions;
source and destination projections to 256; dropout 0.5; one hidden layer; cross-entropy over the event type; 3 epochs
of Adam with learning rate $10^{-3}$ on at most 4 million sampled training events). Node score = maximum event loss.
\item Explanation baselines: KernelSHAP with 20 $k$-means background points and 200 samples; TreeSHAP on the
Isolation Forest; occlusion replaces a channel's features by their training medians; reconstruction-error grouping
takes the channel with the largest mean squared per-feature error.
\item Blast radius: logon weights $\log(1+n)\,e^{-\Delta d/7}$ over earlier days; PageRank restart 0.15, 50
iterations, undirected graph.
\end{itemize}

\subsection{Windows rule prior (process path or command line)}
Each matching rule adds one point to a process; process-creation events add two points to both endpoints when a
document, mail or browser process spawns an interpreter or living-off-the-land binary, and one point when an
interpreter spawns such a tool.
\begin{itemize}\footnotesize
%%WIN_RULES%%
\end{itemize}

\subsection{Linux / FreeBSD / Android rule prior}
Same point scheme; parents are network-facing programs (browsers, mail clients, web, ssh and ftp servers), children
are shells, interpreters and download or discovery tools.
\begin{itemize}\footnotesize
%%LIN_RULES%%
\end{itemize}

\subsection{Ground-truth channels for the explanation study}
An explanation counts as correct if the explained channel is among the ground-truth channels of the host-hour;
host-hours without a mappable channel are excluded. The chance level is the expected accuracy of choosing one of the
four channels uniformly at random, that is, the average of $|\mathrm{GT}\cap\{\text{recon, c2, lateral, exfil}\}|/4$
over the evaluated host-hours.
\begin{itemize}\footnotesize
\item DEDALE (ATT\&CK tactic ID): TA0043, TA0007 $\rightarrow$ recon; TA0011 $\rightarrow$ c2; TA0008 $\rightarrow$
lateral; TA0010 $\rightarrow$ exfil; TA0001, TA0009 $\rightarrow$ none.
\item Unraveled (stage): Reconnaissance $\rightarrow$ recon; Establish Foothold $\rightarrow$ c2; Lateral Movement
$\rightarrow$ lateral; Data Exfiltration $\rightarrow$ exfil; Cover up $\rightarrow$ none.
\item OpTC: flows to a documented attacker server $\rightarrow$ c2; report-timeline entries by keyword (scan
$\rightarrow$ recon; pivot, WMI, RDP, spread $\rightarrow$ lateral; exfiltrat, netcat $\rightarrow$ exfil).
\end{itemize}

\subsection{Values behind Figure~\ref{fig:budget}}
\begin{table}[!htbp]\centering\footnotesize\setlength{\tabcolsep}{3pt}
\caption{Ground-truth nodes found within the top-$k$ alerts (seed 0).}
\begin{tabular}{llcccccc}
\toprule
Dataset & Method & 10 & 25 & 50 & 100 & 250 & 500 \\
\midrule
%%FIG_ROWS%%
\bottomrule
\end{tabular}
\end{table}
